\subsection{OpenMP -- Open Multi-Processing}

OpenMP (Target Offloading)
\begin{itemize}
    \item first OpenMP specification version 1.0 (C/\cpp) October 1998~\cite{openmp_architecture_review_board_openmp_1998}
    \item first OpenMP specification with target offloading version 4.0 July 2013~\cite{openmp_architecture_review_board_openmp_2013}
    \item current OpenMP specification version 6.0 November 2024~\cite{openmp_architecture_review_board_openmp_2024}
    \item Implementations (Auszug)(\url{https://www.openmp.org/resources/openmp-compilers-tools/})
    \begin{itemize}
        \item AMD
        \item Intel
        \item LLVM
        \item GNU
        \item ARM
        \item NVIDIA
        \item PGI
        \item $\dots$
    \end{itemize}
\end{itemize}


\begin{itemize}
    \item \acrfull{openmp}
    \item pragma based
    \item least invasive; pragmas ignored if no OpenMP capable compiler could be found
    \item but comes also with a runtime library 
    \item originally meant for CPUs only
    \item standardized
    \item uses the fork-join model; initial thread spawns worker threads
    \item uses a custom OpemMP thread pool
    \item task- and loop-parallelism (only latter is used)
\end{itemize}


\begin{advantagebox}
    \begin{itemize}
        \item standardized
        \item supports a wide variety of hardware platforms
        \item non-intrusive and easy-to-use
    \end{itemize}
\end{advantagebox}

\begin{disadvantagebox}
    \begin{itemize}
        \item very high-level, i.e., some optimizations not possible
    \end{itemize}
\end{disadvantagebox}
